% ch09b Screening and plasmons, part b (书页 137-145 / 分册p023-031)
%--TAIL--
在小 $\omega$ 下的行为。在我们的 $\varepsilon_{\mathrm{R}}$ 表达式中令 $\omega=0$，得到
\begin{equation*}
\varepsilon(x,\omega=0)=\varepsilon_{\mathrm{R}}(x,\omega=0)=1+\frac{\kappa_{\mathrm{TF}}^{2}}{k^{2}}\left[\frac{1}{2}+\frac{1-x^{2}}{4x}\ln\left|\frac{1+x}{1-x}\right|\right],
\tag{9.123}
\end{equation*}
这就是静态极限下（$x=k/2k_{\mathrm{F}}$）的介电函数。$\varepsilon(x,\omega=0)=1+\kappa_{\mathrm{TF}}^{2}/k^{2}$ 在 $x=0$ 处的极限，正是对原点处正电荷屏蔽的托马斯–费米近似。

与托马斯–费米情形一样，我们可以构造出由这类屏蔽效应产生的空间势。为此，我们利用有效场的方程 $U_{\mathrm{eff}}(\mathbf{k},\omega)=\varepsilon^{-1}(\mathbf{k},\omega)U(\mathbf{k},\omega)$。对电子气，$U(\mathbf{k},\omega=0)=4\pi e^{2}/k^{2}$。若代入我们的 $\varepsilon(\mathbf{k},\omega=0)$ 表达式，便发现
\begin{equation*}
\phi_{\mathrm{eff}}(\mathbf{r})=4\pi e\int\frac{\mathrm{d}\mathbf{k}}{(2\pi)^{3}}\,\frac{\mathrm{e}^{-\mathrm{i}\mathbf{k}\cdot\mathbf{r}}}{k^{2}+\kappa_{\mathrm{TF}}^{2}Q(\mathbf{k})}
\tag{9.124}
\end{equation*}
即有效势的空间依赖关系，其中 $Q(\mathbf{k})$ 表示式 (9.123) 中方括号项的傅里叶变换。当 $k=2k_{\mathrm{F}}$ 时，$Q(\mathbf{k})$ 对数发散。这一发散对静电势给出如下形式的贡献：
\begin{equation*}
\phi_{\mathrm{eff}}\sim\frac{\cos 2k_{\mathrm{F}}r}{r^{3}}
\tag{9.125}
\end{equation*}
（在 $r\to\infty$ 时）。静电势的这种振荡行为纯粹是屏蔽的后果，称为弗里德德尔振荡（Friedel oscillation）[F1954]。于是在长距离处，我们发现电荷并未被屏蔽到足以给出托马斯–费米理论的 $\mathrm{e}^{-\kappa_{\mathrm{TF}}r}/r$ 行为。静电势的代数式衰减意味着：一个局域的外来电荷会影响到电子气中每一处的电荷密度。科恩（Kohn）[K1959] 最先论证，屏蔽静电势的这种缓慢衰减源于费米面的陡峭性。这一效应会在金属的声子谱中显现，条件是激发具有净动量转移 $k>2k_{\mathrm{F}}$。他还与卢廷格（Luttinger）一道 [KL1965] 指出，$\phi_{\mathrm{eff}}$ 的负贡献会在 $T=0$ 的电子气中引起超导失稳。这一观察意义重大，因为它表明：如果任其自然，具有裸排斥相互作用的电子气无需声子的协助也能变为超导。

\section{Kubo公式：电导率}

从线性响应形式体系出发，我们还可以导出由 Kubo [K1957] 给出的、把电导率张量与电流联系起来的关系式。辐射与物质相互作用的一般哈密顿量为
\begin{equation*}
H=\frac{1}{2m}\sum_{i}\left(\mathbf{p}_{i}-\frac{e}{c}\mathbf{A}(\mathbf{r}_{i},t)\right)^{2}+H_{\mathrm{other}},
\tag{9.126}
\end{equation*}
其中包含矢量势 $\mathbf{A}(\mathbf{r}_{i},t)$ 与 $H_{\mathrm{other}}$，后者代表描述多体系统的所有其余各项。假设 $H_{\mathrm{other}}$ 与粒子的动量无关，则速度为
\begin{equation*}
\mathbf{v}_{i}=\nabla_{\mathbf{p}_{i}}H=\frac{1}{m}\left(\mathbf{p}_{i}-\frac{e}{c}\mathbf{A}(\mathbf{r}_{i})\right),
\tag{9.127}
\end{equation*}
因而总电流
\begin{equation*}
\int\mathrm{d}\mathbf{r}\,\mathbf{J}(\mathbf{r},t)=\sum_{i}e\mathbf{v}_{i}=\sum_{i}\frac{e}{m}\left(\mathbf{p}_{i}-\frac{e}{c}\mathbf{A}(\mathbf{r}_{i})\right)
\tag{9.128}
\end{equation*}
包含两个不同的贡献。顺磁项
\begin{equation*}
\mathbf{j}(\mathbf{r},t)=\frac{e}{2m}\sum_{i}\left(\mathbf{p}_{i}\delta(\mathbf{r}-\mathbf{r}_{i})+\delta(\mathbf{r}-\mathbf{r}_{i})\mathbf{p}_{i}\right)
\tag{9.129}
\end{equation*}
通常写成对称化的形式，而抗磁部分
\begin{equation*}
\mathbf{j}_{\mathrm{D}}(\mathbf{r},t)=-\frac{e^{2}}{mcV}\sum_{i}\mathbf{A}(\mathbf{r}_{i},t)
\tag{9.130}
\end{equation*}
则直接正比于矢量势。我们最终感兴趣的是电流的热平均：$\langle\mathbf{J}(\mathbf{r},t)\rangle_{A}$，下标表示计及矢量势的存在。由于我们只关心线性响应理论，可以把抗磁项近似为
\begin{equation*}
\mathbf{j}_{\mathrm{D}}=-\frac{e^{2}n_{\mathrm{e}}}{mc}\mathbf{A}(\mathbf{r},t)
\tag{9.131}
\end{equation*}
以及
\begin{equation*}
\langle\mathbf{J}(\mathbf{r},t)\rangle_{A}=\langle\mathbf{j}(\mathbf{r},t)\rangle_{A}+\mathbf{j}_{\mathrm{D}}.
\tag{9.132}
\end{equation*}

为用电流表述电导率，注意到所有涉及矢量势的项都以如下形式进入哈密顿量：
\begin{equation*}
H'=-\frac{1}{c}\int\mathrm{d}\mathbf{r}\,\mathbf{J}(\mathbf{r},t)\cdot\mathbf{A}(\mathbf{r},t)
\tag{9.133}
\end{equation*}
与密度响应完全类似，把式 (9.52) 应用于电流的期望值，得到
\begin{equation*}
\langle J_{\alpha}(\mathbf{r},t)\rangle_{A}=\frac{\mathrm{i}}{\hbar}\int\mathrm{d}\mathbf{r}'\int_{-\infty}^{\infty}\mathrm{d}t'\,\theta(t-t')\,\langle[j_{\beta}(\mathbf{r}',t'),j_{\alpha}(\mathbf{r},t)]\rangle A_{\beta}(\mathbf{r}',t')-\frac{n_{\mathrm{e}}e^{2}}{m}A_{\alpha}(\mathbf{r},t)
\tag{9.134}
\end{equation*}
作为平均电流的工作表达式。电流与外加电场之间满足线性关系
\begin{equation*}
\langle J_{\alpha}(\mathbf{r},t)\rangle_{A}=\int\mathrm{d}\mathbf{r}'\int_{-\infty}^{\infty}\mathrm{d}t'\,\sigma_{\alpha\beta}(\mathbf{r}-\mathbf{r}',t-t')E_{\beta}(\mathbf{r}',t'),
\tag{9.135}
\end{equation*}
比例常数即电导率张量。由于 $\mathbf{E}(\mathbf{r},t)=-(1/c)\,\partial\mathbf{A}(\mathbf{r},t)/\partial t$，矢量势可以换成 $c\mathbf{E}/\mathrm{i}\omega$。把这一替换代入式 (9.134)，便可辨认出
\begin{align*}
\sigma_{\alpha\beta}(\mathbf{q},\omega)&=\frac{1}{\hbar\omega}\int_{0}^{\infty}\mathrm{d}t\,\mathrm{e}^{-\mathrm{i}\omega t}\,\langle[j_{\beta}(-\mathbf{q},0),j_{\alpha}(\mathbf{q},t)]\rangle-\frac{1}{\mathrm{i}\omega}\frac{n_{\mathrm{e}}e^{2}}{m}\delta_{\alpha\beta}\\
&=\frac{1}{\mathrm{i}\hbar\omega}\chi_{j_{\alpha}j_{\beta}}(\mathbf{q},\omega)-\frac{1}{\mathrm{i}\omega}\frac{n_{\mathrm{e}}e^{2}}{m}\delta_{\alpha\beta}
\tag{9.136}
\end{align*}
为电导率。这里 $\chi_{j_{\alpha}j_{\beta}}(\mathbf{q},\omega)$ 是电流–电流响应函数。电导率的实部为
\begin{equation*}
\mathrm{Re}\,\sigma_{\alpha\beta}(\mathbf{q},\omega)=\frac{1}{\hbar\omega}\,\mathrm{Im}\,\chi_{j_{\alpha}j_{\beta}}(\mathbf{q},\omega).
\tag{9.137}
\end{equation*}
只有在超导体中，抗磁项才得以保留，且电导率的虚部发散。因此，在金属中，式 (9.136) 右端第一项给出的虚部贡献恰好与抗磁响应相消。直流极限
\begin{equation*}
\sigma_{\alpha\beta}^{\mathrm{dc}}=\lim_{\omega\to0}\lim_{q\to0}\left(\frac{1}{\hbar\omega}\,\mathrm{Im}\,\chi_{j_{\alpha}j_{\beta}}(\mathbf{q},\omega)\right),
\tag{9.138}
\end{equation*}
通过先取 $q\to0$ 得到，此时体系成为空间均匀的状态。

电导率的另一种推导方式是把微扰
\begin{equation*}
H'=-e\int\mathbf{r}\cdot E(\mathbf{r},t)\,\mathrm{d}\mathbf{r},
\tag{9.139}
\end{equation*}
直接用外加电场表出。借助式 (9.134) 与线性响应的关键结果式 (9.50)，我们辨认出
\begin{equation*}
\sigma_{\alpha\beta}(t,t')=-\frac{e}{\mathrm{i}\hbar}\langle[J_{\alpha}(t),r_{\beta}(t')]\rangle
\tag{9.140}
\end{equation*}
为电导率。电流就是
\begin{equation*}
\mathbf{J}=\frac{e}{m}\dot{\mathbf{r}},
\tag{9.141}
\end{equation*}
其中 $\dot{\mathbf{r}}=\mathbf{p}-e\mathbf{A}/c$。电导率的这一形式对于推导所谓的 f 求和规则（f-sum rule）特别有用。先写出电导率的傅里叶变换
\begin{equation*}
\sigma_{\alpha\beta}(\omega)=\int_{-\infty}^{\infty}\mathrm{d}(t-t')\,\mathrm{e}^{\mathrm{i}\omega(t-t')}\sigma_{\alpha\beta}(t-t'),
\tag{9.142}
\end{equation*}
于是，电导率的积分权重
\begin{align*}
\frac{1}{2\pi}\int_{-\infty}^{\infty}\sigma_{\alpha\beta}(\omega)\mathrm{d}\omega&=\frac{1}{2\pi}\int_{-\infty}^{\infty}\mathrm{d}(t-t')\,2\pi\delta(t-t')\,\sigma_{\alpha\beta}(t,t')\\
&=-\frac{1}{\mathrm{i}\hbar}\frac{e^{2}}{m}\langle[\dot{r}_{\alpha},r_{\beta}]\rangle\\
&=\delta_{\alpha\beta}\frac{n_{\mathrm{e}}e^{2}}{m}
\tag{9.143}
\end{align*}
满足一个求和规则。使用这一求和规则时，我们显式地假定动能取连续形式 $p^{2}/2m$，由此才有 $\dot{\mathbf{r}}=\mathbf{p}-e\mathbf{A}/c$。

\section{等离体的阻止本领}

把一个具有一定入射能量的电子注入等离体中时，由于它与等离体中电子的库仑相互作用，预期它会被减速。基于这些理由，贝特（Bethe）论证道，注入电子的能量损失率应当正比于 $|V_{\mathrm{int}}(\mathbf{p})|^{2}$，其中 $V_{\mathrm{int}}$ 是等离体与电子之间的库仑相互作用。对库仑势，如前所示，$V_{\mathrm{int}}\sim p^{-2}$。对所有入射动量值求和，我们发现能量损失为
\begin{equation*}
\frac{\mathrm{d}E}{\mathrm{d}t}\propto\int\mathrm{d}\mathbf{p}\,p\,|V_{\mathrm{int}}(\mathbf{p})|^{2}\propto\int_{0}^{p_{\mathrm{F}}}\frac{\mathrm{d}p}{p}.
\tag{9.144}
\end{equation*}
该积分在下限处对数发散。这样一来，这一简单论证便给出发散的能量损失，这显然不对。我们立即看到，对于屏蔽相互作用，下限处的发散将会消失，从而使 $\mathrm{d}E/\mathrm{d}t$ 有限。这就是贝特方法的主要失败之处。下面我们以一种严格的方式表述该问题，通过把屏蔽效应包括进来而绕开贝特发散。

手头的物理问题是：处于某个初态 $|i\rangle$ 的金属，与一个具有初始动量 $\mathbf{p}$ 的电子发生碰撞。与金属相互作用之后，电子将具有新动量 $\mathbf{p}-\hbar\mathbf{k}$，而金属将处于某个新态 $|f\rangle$。我们假定金属的态构成一组完备的正交归一集合 $\langle n|m\rangle=\delta_{nm}$。在费米黄金定则（Fermi's golden rule）的层面上，初末态之间的跃迁率为
\begin{equation*}
W_{\mathbf{p},\mathbf{p}-\hbar\mathbf{k}}=\frac{2\pi}{\hbar}\sum_{f}\left|\langle f,\mathbf{p}-\hbar\mathbf{k}|V_{\mathrm{int}}|i,\mathbf{p}\rangle\right|^{2}\delta\left(\frac{\mathbf{p}^{2}}{2m}+E_{i}-\frac{(\mathbf{p}-\hbar\mathbf{k})^{2}}{2m}-E_{f}\right)
\end{equation*}
其中 $E_{i}$ 与 $E_{f}$ 分别是态 $|i\rangle$ 与 $|f\rangle$ 的总能量。相互作用能
\begin{equation*}
V_{\mathrm{int}}(\mathbf{r})=\int\mathrm{d}\mathbf{r}'\,\left[n(\mathbf{r}')-n_{\mathrm{e}}\right]\frac{e^{2}}{|\mathbf{r}'-\mathbf{r}|},
\tag{9.145}
\end{equation*}
既包括离子的也包括电子的库仑能。电子的初态是形如 $|\mathbf{p}\rangle=\mathrm{e}^{\mathrm{i}\mathbf{p}\cdot\mathbf{r}/\hbar}/\sqrt{V}$ 的平面波，末态则为 $|\mathbf{p}-\hbar\mathbf{k}\rangle=\mathrm{e}^{\mathrm{i}(\mathbf{p}-\hbar\mathbf{k})\cdot\mathbf{r}/\hbar}/\sqrt{V}$。利用这些态，我们把式 (9.145) 中的矩阵元改写为
\begin{align*}
\langle f,\mathbf{p}-\hbar\mathbf{k}|V_{\mathrm{int}}|i,\mathbf{p}\rangle&=\int\mathrm{d}\mathbf{r}\,\mathrm{d}\mathbf{r}'\,\langle f|n(\mathbf{r}')-n_{\mathrm{e}}|i\rangle\,\frac{\mathrm{e}^{\mathrm{i}\mathbf{k}\cdot\mathbf{r}}}{V}\,\frac{e^{2}}{|\mathbf{r}-\mathbf{r}'|}\\
&=\frac{4\pi e^{2}}{k^{2}}\int\frac{\mathrm{d}\mathbf{r}}{V}\,\mathrm{e}^{\mathrm{i}\mathbf{k}\cdot\mathbf{r}}\langle f|n(\mathbf{r})|i\rangle.
\tag{9.146}
\end{align*}
利用 $\delta$ 函数的积分表示
\begin{equation*}
2\pi\hbar\delta(\hbar\omega+E_{i}-E_{f})=\int_{-\infty}^{\infty}\mathrm{e}^{\mathrm{i}(\hbar\omega+E_{i}-E_{f})t/\hbar}\,\mathrm{d}t,
\end{equation*}
我们把跃迁率改写为
\begin{align*}
W_{\mathbf{p},\mathbf{p}-\hbar\mathbf{k}}&=\left(\frac{4\pi e^{2}}{V\hbar k^{2}}\right)^{2}\int_{-\infty}^{\infty}\mathrm{e}^{\mathrm{i}(\hbar\omega+E_{i}-E_{f})t/\hbar}\mathrm{d}t\sum_{f\neq i}\left|\langle f|\int\mathrm{d}\mathbf{r}\,n(\mathbf{r})\mathrm{e}^{\mathrm{i}\mathbf{k}\cdot\mathbf{r}}|i\rangle\right|^{2}\\
&=\left(\frac{4\pi e^{2}}{V\hbar k^{2}}\right)^{2}\int_{-\infty}^{\infty}\mathrm{d}t\,\mathrm{e}^{\mathrm{i}\omega t}\sum_{f\neq i}\langle f|n(\mathbf{k})|i\rangle\langle i|\mathrm{e}^{\mathrm{i}E_{i}\frac{t}{\hbar}}n(\mathbf{k})\mathrm{e}^{-\mathrm{i}E_{f}\frac{t}{\hbar}}|f\rangle
\tag{9.147}
\end{align*}
其中 $\hbar\omega=\epsilon_{\mathbf{p}}-\epsilon_{\mathbf{p}-\hbar\mathbf{k}}$。在相互作用表象中，
\begin{equation*}
\widehat{n}(\mathbf{k},t)=\mathrm{e}^{\mathrm{i}H_{0}\frac{t}{\hbar}}n(\mathbf{k})\mathrm{e}^{-\mathrm{i}H_{0}\frac{t}{\hbar}}.
\tag{9.148}
\end{equation*}
于是
\begin{equation*}
W_{\mathbf{p},\mathbf{p}-\hbar\mathbf{k}}=\left(\frac{4\pi e^{2}}{V\hbar k^{2}}\right)^{2}\int_{-\infty}^{\infty}\mathrm{e}^{\mathrm{i}\omega t}\,\mathrm{d}t\sum_{f\neq i}\langle i|\widehat{n}(\mathbf{k},t)|f\rangle\langle f|n(\mathbf{k})|i\rangle,
\tag{9.149}
\end{equation*}
利用结构函数的定义以及金属态的完备性关系 $\sum|f\rangle\langle f|=1$，它可以简化为
\begin{equation*}
W_{\mathbf{p},\mathbf{p}-\hbar\mathbf{k}}=\left(\frac{4\pi e^{2}}{\hbar k^{2}}\right)^{2}\frac{n_{\mathrm{e}}}{V}\left[S(\mathbf{k},\omega)-\frac{1}{V}2\pi\delta(\omega)n_{\mathrm{e}}\right]
\tag{9.150}
\end{equation*}
正如所料，决定体系对入射电子的响应的，正是动力学结构因子。屏蔽效应已隐含地包括在 $S(\mathbf{k},\omega)$ 之中。贝特结果来自跃迁率的零频部分。

我们最感兴趣的是电子交给等离体的能量损失率。它由对所有能量差 $\epsilon_{\mathbf{p}}-\epsilon_{\mathbf{p}-\hbar\mathbf{k}}$ 按跃迁率 $W$ 加权求和来确定：
\begin{align*}
\frac{\mathrm{d}E}{\mathrm{d}t}&=-\sum_{\mathbf{k}}(\epsilon_{\mathbf{p}}-\epsilon_{\mathbf{p}-\hbar\mathbf{k}})W_{\mathbf{p},\mathbf{p}-\hbar\mathbf{k}}\\
&=-\frac{\hbar n_{\mathrm{e}}}{V}\sum_{\mathbf{k}}\omega\left(\frac{4\pi e^{2}}{\hbar k^{2}}\right)^{2}S(\mathbf{k},\omega)\\
&=-n_{\mathrm{e}}\int\frac{\mathrm{d}\mathbf{k}}{(2\pi)^{3}}\left(\frac{4\pi e^{2}}{\hbar k^{2}}\right)^{2}\int_{-\infty}^{\infty}\omega\,\mathrm{d}\omega\,S(\mathbf{k},\omega)\,\delta(\hbar\omega-\epsilon_{\mathbf{p}}+\epsilon_{\mathbf{p}-\hbar\mathbf{k}}).
\tag{9.151}
\end{align*}
为求这个量，我们改用极坐标，并先对 $\mathbf{p}$ 与 $\mathbf{p}-\hbar\mathbf{k}$ 之间的夹角 $\theta$ 积分：
\begin{align*}
2\pi\int_{0}^{\pi}\mathrm{d}\theta\,\sin\theta\,\delta\left(\hbar\omega-\epsilon_{\mathbf{p}}+\epsilon_{\mathbf{p}-\hbar\mathbf{k}}\right)&=2\pi\int_{-1}^{1}\mathrm{d}x\,\delta\left(\hbar\omega+\frac{\hbar^{2}k^{2}}{2m}-\frac{p\hbar kx}{m}\right)\\
&=\frac{2\pi m}{p\hbar k}\Theta\left(kv_{\mathrm{p}}-\left|\omega+\hbar k^{2}/2m\right|\right).
\tag{9.152}
\end{align*}
利用这一结果，交给等离体的能量损失率简化为
\begin{equation*}
\frac{1}{v_{\mathrm{p}}}\frac{\mathrm{d}E}{\mathrm{d}t}=\frac{-4e^{4}n_{\mathrm{e}}}{\hbar^{2}v_{\mathrm{p}}^{2}}\int_{0}^{\infty}\frac{\mathrm{d}k}{k^{3}}\int_{-kv_{\mathrm{p}}-\frac{\hbar k^{2}}{2m}}^{kv_{\mathrm{p}}-\frac{\hbar k^{2}}{2m}}\mathrm{d}\omega\,\omega S(\mathbf{k},\omega),
\tag{9.153}
\end{equation*}
其中 $v_{\mathrm{p}}$ 是入射电子的速度。

电子被等离体完全阻止，最可能发生在电子气作集体行为之时，也就是说，在等离体振荡占主导之时。因此，当 $|\omega|=\omega_{\mathrm{p}}$（等离体频率）时，电子气获得最大的阻止本领（stopping power）。于是我们要寻求结构函数在高频极限下的表达式。从介电函数的定义出发（见式 (9.75) 与 (9.93)），我们把 $\varepsilon(\mathbf{k},\omega)^{-1}$ 的虚部
\begin{equation*}
\mathrm{Im}\,\varepsilon^{-1}=-\frac{2\pi e^{2}n_{\mathrm{e}}}{\hbar k^{2}}\left(1-\mathrm{e}^{-\beta\hbar\omega}\right)S(\mathbf{k},\omega),
\tag{9.154}
\end{equation*}
用结构函数表出，并对介电函数使用高频展开
\begin{equation*}
\varepsilon(\omega)\sim\lim_{\eta\to0}\left(1-\frac{\omega_{\mathrm{p}}^{2}}{(\omega+\mathrm{i}\eta)^{2}}\right).
\tag{9.155}
\end{equation*}
在此极限下，$\varepsilon^{-1}$ 的虚部
\begin{align*}
\mathrm{Im}\,\varepsilon^{-1}(\omega)&=\lim_{\eta\to0}\,\mathrm{Im}\left[\frac{\omega^{2}}{(\omega+\mathrm{i}\eta)^{2}-\omega_{\mathrm{p}}^{2}}\right]\\
&=-\frac{\pi\omega_{\mathrm{p}}}{2}\left[\delta(\omega-\omega_{\mathrm{p}})-\delta(\omega+\omega_{\mathrm{p}})\right],
\tag{9.156}
\end{align*}
是位于 $\pm\omega_{\mathrm{p}}$ 处的两个 $\delta$ 函数之和。借助式 (9.154)，立即可见在 $k\to0$ 极限下
\begin{equation*}
S(\mathbf{k},\omega)=\frac{\hbar\pi k^{2}}{m}\,\frac{\left[\delta(\omega-\omega_{\mathrm{p}})-\delta(\omega+\omega_{\mathrm{p}})\right]}{1-\mathrm{e}^{-\beta\hbar\omega}}.
\tag{9.157}
\end{equation*}
式 (9.153) 中的 $\omega$ 积分如今直截了当：
\begin{equation*}
\int_{\omega_{\ell}}^{\omega_{\mathrm{u}}}\mathrm{d}\omega\,\omega S(\mathbf{k},\omega)=\frac{\hbar\pi k^{2}}{m}(1+g_{\mathrm{p}})\,\Theta(\omega_{\ell}<\omega_{\mathrm{p}}<\omega_{\mathrm{u}})-\hbar\frac{\pi k^{2}}{m}g_{\mathrm{p}}\,\Theta(\omega_{\ell}<-\omega_{\mathrm{p}}<\omega_{\mathrm{u}}),
\tag{9.158}
\end{equation*}
其中 $\omega_{\ell}=-kv_{\mathrm{p}}-\hbar k^{2}/2m$，$\omega_{\mathrm{u}}=kv_{\mathrm{p}}-\hbar k^{2}/2m$。在计算这一积分时，我们引入了 $g_{\mathrm{p}}=(\mathrm{e}^{\beta\hbar\omega_{\mathrm{p}}}-1)^{-1}$，它决定在温度 $T$ 下被热激发的等离激元数目。能量损失变换为
\begin{align*}
\frac{1}{v_{\mathrm{p}}}\frac{\mathrm{d}E}{\mathrm{d}t}&=-\left(\frac{\omega_{\mathrm{p}}e}{\hbar v_{\mathrm{p}}}\right)^{2}\int\frac{\mathrm{d}k}{k}\,(1+g_{\mathrm{p}})\,\Theta\left(\hbar kv_{\mathrm{p}}-\frac{(\hbar k)^{2}}{2m}-\hbar\omega_{\mathrm{p}}>0\right)\\
&\quad-g_{\mathrm{p}}\,\Theta\left(\frac{(\hbar k)^{2}}{2m}-\hbar kv_{\mathrm{p}}<\hbar\omega_{\mathrm{p}}<\frac{(\hbar k)^{2}}{2m}+\hbar kv_{\mathrm{p}}\right).
\tag{9.159}
\end{align*}
第一项代表产生一个等离激元时的能量损失，第二项则是介质中热激发的等离激元转移给电子的能量。当 $T\to0$ 时，等离激元被热激发的概率趋于零，因此等离激元只能由入射电子来激发。在此极限下，能量损失取如下简单形式：
\begin{align*}
\frac{1}{v_{\mathrm{p}}}\frac{\mathrm{d}E}{\mathrm{d}t}&=-\left(\frac{\omega_{\mathrm{p}}e}{\hbar v_{\mathrm{p}}}\right)^{2}\int\frac{\mathrm{d}k}{k}\,\Theta\left(\hbar kv_{\mathrm{p}}-\frac{\hbar^{2}k^{2}}{2m}-\hbar\omega_{\mathrm{p}}>0\right)\\
&=-\left(\frac{\omega_{\mathrm{p}}e}{\hbar v_{\mathrm{p}}}\right)^{2}\ln\frac{k_{+}}{k_{-}},
\tag{9.160}
\end{align*}
其中 $k_{\pm}$ 是
\begin{equation*}
\hbar^{2}k^{2}-2m\hbar kv_{\mathrm{p}}+2m\hbar\omega_{\mathrm{p}}=0,
\tag{9.161}
\end{equation*}
的解，或者等价地，
\begin{equation*}
\hbar k_{\pm}=p\pm\sqrt{p^{2}-2m\hbar\omega_{\mathrm{p}}}.
\tag{9.162}
\end{equation*}
要使入射电子能够激发等离激元，需要 $p^{2}/2m>\hbar\omega_{\mathrm{p}}$\,\footnote{译注：原书此处印作 $p^{2}/2m>\omega_{\mathrm{p}}$，按量纲应为 $\hbar\omega_{\mathrm{p}}$。}。在此极限下展开 $k_{\pm}$，得到
\begin{equation*}
k_{\pm}=p\pm(p-m\hbar\omega_{\mathrm{p}}/p)=\left\{\begin{array}{c}2p-m\hbar\omega_{\mathrm{p}}/p\\ m\hbar\omega_{\mathrm{p}}/p\end{array}\right\}.
\tag{9.163}
\end{equation*}
若 $k\lesssim a^{-1}$（$a$ 为粒子间距），我们预期集体性的等离激元振荡将不复存在。于是应把 $k_{+}$ 在 $2p_{\mathrm{F}}$ 处截断。由此得到
\begin{equation*}
\frac{\mathrm{d}E}{\mathrm{d}t}=-\frac{\omega_{\mathrm{p}}^{2}e^{2}}{v_{\mathrm{p}}\hbar^{2}}\ln\frac{2pp_{\mathrm{F}}}{m\hbar\omega_{\mathrm{p}}},
\tag{9.164}
\end{equation*}
它完全良态且有限，与贝特结果不同。

\section*{小结}
\addcontentsline{toc}{section}{小结}

电子气展现出形形色色的集体长波现象，例如等离体振荡与屏蔽，它们显著地改变了电子气的独立电子图像。就屏蔽而言，电子气作集体行为，把裸相互作用从 $V_{\mathbf{q}}$ 降低为
\begin{equation*}
V_{\mathrm{eff}}=\frac{V_{\mathbf{q}}}{\varepsilon(q,\omega)},
\tag{9.165}
\end{equation*}
其中 $\varepsilon(q,\omega)$ 是介电函数。在金属密度下，RPA 足以描述屏蔽效应。在这一近似中，把无关联的粒子–空穴散射按微扰论的所有阶求和，得到极化函数，或称极化气泡（polarization bubble）。RPA 不包括交换效应。在足够稀释的电子体系（$r_{\mathrm{s}}>3$）中，库仑能与电子动能相当甚至更大，此时需要超出 RPA 的理论。诸如局域场方法（local-field approach）[STLS1968] 这类试图对每个电子周围的关联空穴建模的方法，在这一情形中已被证明相当成功（见习题 9.5）。

\section*{习题}
\addcontentsline{toc}{section}{习题}

\textbf{习题 9.1}\quad 对位于原点的 $+Q$ 电荷，在 $d=1$ 和 $d=2$ 维中重做托马斯–费米屏蔽问题。假设未屏蔽势为库仑势。显式计算 $d=1$ 和 $d=2$ 中的托马斯–费米屏蔽长度。显式证明在 $d=2$ 中，托马斯–费米屏蔽长度与密度无关，且由 $1/\kappa_{\mathrm{TF}}=a_{0}/2$ 给出。

\textbf{习题 9.2}\quad 在屏蔽的托马斯–费米处理中，显式计算式 (9.7) 中正比于 $U_{\mathrm{eff}}^{2}$ 的项。在什么条件下可以忽略这一项？

\textbf{习题 9.3}\quad 用等离激元一节中发展的方法证明，$d=2$ 中的等离激元由按 $\omega_{\mathrm{p}}\propto\sqrt{k}$ 标度的色散关系描述。

\textbf{习题 9.4}\quad 推导式 (9.31) 与式 (9.32)，并利用这些结果推导式 (9.33)。\footnote{译注：原文作 Derive Eqs.~(9.31)，未列出式 (9.32)；按句中“这些结果”的复数语气，当指式 (9.31) 与 (9.32)。}

\textbf{习题 9.5}\quad 在 $d=1$ 和 $d=2$ 中对等离体的阻止本领重做贝特的论证。为什么结果是有限的？

\textbf{习题 9.6}\quad 在小波矢极限下，介电函数取如下形式 [N1964]
\begin{equation*}
\lim_{q\to0}\varepsilon(q,0)=1+\frac{\kappa_{\mathrm{TF}}^{2}}{q^{2}}\,\frac{\alpha}{\alpha_{\mathrm{F}}},
\tag{9.166}
\end{equation*}
其中 $\alpha$ 是电子气的压缩率，$\alpha_{\mathrm{F}}$ 是自由体系的压缩率。\footnote{译注：原书此处印作 $\kappa_{\mathrm{F}}$，与式 (9.166) 及下文的 $\alpha_{\mathrm{F}}$ 不一致，按本意统一作 $\alpha_{\mathrm{F}}$。}对自由体系，$\alpha_{\mathrm{F}}=3/(n_{\mathrm{e}}\epsilon_{\mathrm{F}})$。证明在 RPA 下 $\alpha/\alpha_{\mathrm{F}}=1$。哈伯德 [H1957] 表明，交换空穴效应可以通过在介电函数中引入因子 $G(q)=q^{2}/2(q^{2}+k_{\mathrm{F}}^{2})$ 来包括：
\begin{equation*}
\varepsilon(q,\omega)=1-\frac{v_{\mathbf{q}}\chi_{nn}^{0}}{1+v_{\mathbf{q}}\chi_{nn}^{0}G(q)}.
\tag{9.167}
\end{equation*}
利用式 (9.166) 与 (9.167)，推导哈伯德近似下压缩率的表达式。画出 $\alpha$ 随 $r_{\mathrm{s}}$ 变化的函数图像。确定 $\alpha$ 改变符号处的 $r_{\mathrm{s}}$ 值。物理上，负压缩率 $\alpha<0$ 意味着什么？

\section*{参考文献}
\addcontentsline{toc}{section}{参考文献}

\begin{itemize}[leftmargin=2em,itemsep=0pt,parsep=0pt,label={}]
\item [BP1953] D. Bohm and D. Pines, \emph{Phys. Rev.} \textbf{92}, 609 (1953).
\item [F1928] E. Fermi, \emph{Z. Phys.} \textbf{48}, 73 (1928).
\item [F1954] H. Friedel, \emph{Adv. Phys.} \textbf{3}, 446 (1954).
\item [H1957] J. Hubbard, \emph{Proc. R. Soc. London Ser. A} \textbf{243}, 336 (1957).
\item [K1959] W. Kohn, \emph{Phys. Rev. Lett.} \textbf{2}, 393 (1959).
\item [KL1965] W. Kohn and J. H. Luttinger, \emph{Phys. Rev. Lett.} \textbf{15}, 524 (1965).
\item [K1957] R. Kubo, \emph{J. Phys. Soc. Japan} \textbf{12}, 570 (1957).
\item [L1954] J. Lindhard, \emph{Kgl. Danske Videnskab. Selskab, Mat-fys. Medd.} \textbf{28}, 8 (1954).
\item [N1964] P. Nozi\`eres, \emph{Interacting Fermi Systems} (Benjamin, Reading, MA., 1964), p.~287.
\item [P1963] D. Pines, \emph{Elementary Excitations in Solids} (Benjamin, Reading, MA., 1963).
\item [STLS1968] K. S. Singwi, M. P. Tosi, R. H. Rand, and A. Sj\"olander, \emph{Phys. Rev. Lett.} \textbf{176}, 589 (1968).
\item [T1927] L. H. Thomas, \emph{Proc. Camb. Phil. Soc.} \textbf{23}, 542 (1927).
\end{itemize}
